% chapters/results.tex  — Sections 7, 8, 9, 10, 11

\chapter{Results and Observations}{\label{ch:results}}

\section{System Outputs}

The system successfully provides a unified interface for managing academic operations. Key outputs
include:

\begin{itemize}
    \item \textbf{Role Based Dashboards:} Personalized views displaying schedules, attendance, and pending
    tasks.
    \item \textbf{Attendance Records:} Automated session wise and student wise tracking, including
    participation scores.
    \item \textbf{Feedback Data:} Structured responses aggregated into summaries and analytics.
    \item \textbf{Reports and Visualizations:} Data presented through graphs, tables, and downloadable
    reports.
\end{itemize}

\section{Observations}

Based on system behaviour and usage, several key observations were made:

\begin{itemize}
    \item \textbf{Reduction in Manual Effort:} Automation of attendance tracking and feedback collection
    significantly reduces administrative workload.
    \item \textbf{Improved Data Availability:} Real time access to scheduling, attendance, and feedback data
    improves transparency for all stakeholders.
    \item \textbf{Enhanced Student Engagement Tracking:} The combination of attendance and feedback data
    provides a more comprehensive view of student participation.
    \item \textbf{Streamlined Workflows:} Integration between modules ensures that processes such as
    attendance computation and feedback rollout occur without manual intervention.
    \item \textbf{Centralized Data Management:} Consolidation of academic data into a single system
    simplifies monitoring and reporting.
\end{itemize}

These observations indicate that the system effectively addresses many of the inefficiencies present
in traditional approaches.

\section{Limitations}

Despite its effectiveness, the system has certain limitations:

\begin{itemize}
    \item \textbf{Dependence on Network Infrastructure:} The attendance system relies on WiFi connectivity
    and accurate device detection, which may be affected by network conditions.
    \item \textbf{Device Dependency:} Attendance accuracy depends on students connecting their registered
    devices to the network.
    \item \textbf{Potential Security Concerns:} Although mitigated through verification, MAC based systems
    can be susceptible to spoofing.
    \item \textbf{Limited Scale Testing:} The system has primarily been tested within the CiPD environment
    and may require further validation for large scale deployment.
\end{itemize}

These limitations highlight areas where further refinement and enhancement may be required.

% ─────────────────────────────────────────────────────────────────────────────
\chapter{Challenges and Solutions}
% ─────────────────────────────────────────────────────────────────────────────

\section{Attendance Reliability and Accuracy}

\textbf{Challenge:}

Ensuring reliable attendance detection using WiFi signals was a major challenge. Factors such as
device disconnections, weak signal strength, and intermittent connectivity could lead to inaccurate
presence detection.

\textbf{Solution:}

Instead of relying on a single detection point, a multi interval detection strategy was implemented.
Attendance is computed based on repeated detections across the session duration, reducing the
impact of temporary disconnections. Additionally, signal strength filtering and threshold based logic
were used to improve detection accuracy.

\section{Network Access and Infrastructure Constraints}

\textbf{Challenge:}

Initial attempts to use the institute's existing WiFi network for attendance tracking were restricted
due to permission limitations. Continuous monitoring of connected devices was not allowed.

\textbf{Solution:}

A dedicated WiFi network was established specifically for the CiPD program. A separate router was
configured to provide controlled access, enabling reliable detection of connected devices without
violating network policies. This approach provided full control over the attendance data pipeline.

\section{Hardware and Firmware Limitations}

\textbf{Challenge:}

To enhance network level data collection, attempts were made to install custom firmware
(OpenWRT) on the router. However, the specific hardware used did not support this modification.

\textbf{Solution:}

The system was adapted to work within the constraints of the existing router firmware. Instead of
deep packet level analysis, the implementation relied on available device detection capabilities,
ensuring functionality without requiring unsupported configurations.

\section{Preventing Proxy Attendance}

\textbf{Challenge:}

A key concern in automated attendance systems is the possibility of proxy attendance, especially
when relying on device based identification.

\textbf{Solution:}

A combination of safeguards was implemented, including:

\begin{itemize}
    \item MAC address registration with administrative verification
    \item Monitoring of detection patterns for anomalies
    \item Provision for manual checks and penalties
\end{itemize}

These measures reduce the likelihood of misuse while maintaining system usability.

\section{Feedback Participation and Engagement}

\textbf{Challenge:}

Ensuring consistent student participation in feedback submission is often difficult, leading to
incomplete or biased data.

\textbf{Solution:}

The feedback system was designed to be lightweight and quick to complete. A notification and
reminder mechanism was implemented to prompt students to submit feedback within a defined
time window, improving participation rates.

\section{Balancing Automation and Control}

\textbf{Challenge:}

While automation reduces manual effort, it also introduces risks if processes fail or produce incorrect
results without oversight.

\textbf{Solution:}

Administrative controls such as manual overrides, monitoring dashboards, and configurable
parameters were incorporated. This ensures that the system remains flexible and allows human
intervention when necessary.

% ─────────────────────────────────────────────────────────────────────────────
\chapter{Future Scope}
% ─────────────────────────────────────────────────────────────────────────────

While the CIPD 360 ERP system provides a functional and integrated solution for academic
management, there are several opportunities for further enhancement and expansion.

One potential improvement is the incorporation of advanced analytics and machine learning
techniques to derive deeper insights from attendance and feedback data. Predictive models could be
used to identify at risk students, analyze engagement trends, and support data driven academic
interventions.

The attendance system can be further strengthened by exploring hybrid approaches, such as
combining WiFi based detection with technologies like Bluetooth or facial recognition. This would
improve accuracy and reduce dependency on a single method. But a more robust approach would be
to integrate this software over to mobile S and utilize the institution network for a full proof
mechanism.

Another area of expansion is the development of a dedicated mobile application, which would
enhance accessibility and enable real time notifications, making the system more convenient for
daily use.

The platform can also be extended to support full scale Learning Management System (LMS)
capabilities, including advanced assignment workflows, peer evaluations, and incorporation of
agentic AI to automate workflows.

From a deployment perspective, the system can be scaled beyond the CiPD program to support
multiple courses, departments, or institutions, with configurable modules to adapt to different
academic structures. Meaning it can be cross institutionalized and can be used to support similar
programs in other colleges.

% ─────────────────────────────────────────────────────────────────────────────
\chapter{Conclusion}
% ─────────────────────────────────────────────────────────────────────────────

This project presented the design and development of CIPD 360 ERP, an integrated academic
management system tailored for the CiPD program at IIIT Delhi. The system was developed to
address key inefficiencies in traditional academic workflows by introducing automation,
centralization, and data driven insights.

By integrating scheduling, attendance tracking, feedback collection, and reporting into a unified
platform, the system successfully streamlines core academic processes. A significant contribution of
this work is the implementation of a WiFi based automated attendance mechanism, which
demonstrates a practical and scalable alternative to conventional attendance systems. In addition,
the incorporation of a structured feedback system and role based dashboards enables continuous
monitoring and evaluation of academic activities.

The project also highlights the importance of iterative development and user centric design. Through
continuous feedback from stakeholders and multiple rounds of experimentation, the system was
refined to align with real world requirements and constraints. Challenges related to network
limitations, reliability, and system integration were addressed through thoughtful design decisions
and practical trade offs.

Overall, the CIPD 360 ERP system demonstrates the feasibility of building a scalable and modular
academic ERP solution for dynamic educational environments. While certain limitations remain, the
system establishes a strong foundation for future enhancements and broader deployment.

% ─────────────────────────────────────────────────────────────────────────────
\appendix
\chapter{API Endpoints}
% ─────────────────────────────────────────────────────────────────────────────

\section{Authentication APIs}

\begin{table}[!ht]
    \centering
    \begin{tabular}{|l|l|p{6.5cm}|}
    \hline
    \textbf{Method} & \textbf{Endpoint} & \textbf{Description} \\ \hline
    POST  & \texttt{/api/auth/login}              & User login (email/enrollment + password) \\
    POST  & \texttt{/api/auth/signup}             & Student registration \\
    GET   & \texttt{/api/auth/me}                & Get current user profile \\
    POST  & \texttt{/api/auth/logout}            & Logout user \\
    POST  & \texttt{/api/auth/update-password}   & Change password \\
    GET   & \texttt{/api/auth/google/connect}    & Initiate Google OAuth \\
    GET   & \texttt{/api/auth/google/callback}   & Handle OAuth callback \\ \hline
    \end{tabular}
    \caption{Authentication API endpoints}
    \label{tab:api_auth}
\end{table}

\section{User Management APIs}

\begin{table}[!ht]
    \centering
    \begin{tabular}{|l|l|p{6.5cm}|}
    \hline
    \textbf{Method} & \textbf{Endpoint} & \textbf{Description} \\ \hline
    GET   & \texttt{/api/students/profile}  & Fetch student profile \\
    GET   & \texttt{/api/students/settings} & Get user preferences \\
    PATCH & \texttt{/api/students/settings} & Update preferences \\
    PATCH & \texttt{/api/students/mac}      & Register MAC address \\ \hline
    \end{tabular}
    \caption{User management API endpoints}
    \label{tab:api_user}
\end{table}

\section{Scheduling \& Session APIs}

\begin{table}[!ht]
    \centering
    \begin{tabular}{|l|l|p{5.5cm}|}
    \hline
    \textbf{Method} & \textbf{Endpoint} & \textbf{Description} \\ \hline
    \multicolumn{3}{|l|}{\textit{Student}} \\ \hline
    GET & \texttt{/api/calendar/sessions}          & Fetch sessions by date range \\
    GET & \texttt{/api/students/schedule/today}    & Today's schedule \\
    GET & \texttt{/api/students/schedule/week}     & Weekly schedule \\ \hline
    \multicolumn{3}{|l|}{\textit{Admin}} \\ \hline
    GET   & \texttt{/api/admin/schedule}                  & Get all sessions (filtered) \\
    GET   & \texttt{/api/admin/sessions}                  & List sessions \\
    POST  & \texttt{/api/admin/sessions}                  & Create session \\
    PATCH & \texttt{/api/admin/sessions/[id]}             & Update session / mark completed \\
    GET   & \texttt{/api/admin/session-types}             & Fetch session types \\
    POST  & \texttt{/api/admin/session-types}             & Create session type \\
    GET   & \texttt{/api/admin/lookup}                    & Fetch courses, faculty, venues, etc. \\
    GET   & \texttt{/api/admin/skills}                    & Get skills \\
    POST  & \texttt{/api/admin/skills}                    & Create skill \\
    GET   & \texttt{/api/admin/categories}                & Get categories \\
    POST  & \texttt{/api/admin/categories}                & Create category \\
    POST  & \texttt{/api/admin/sessions/rollout-feedback} & Trigger feedback rollout \\ \hline
    \end{tabular}
    \caption{Scheduling and session API endpoints}
    \label{tab:api_schedule}
\end{table}

\section{Attendance APIs}

\begin{table}[!ht]
    \centering
    \begin{tabular}{|l|l|p{5.5cm}|}
    \hline
    \textbf{Method} & \textbf{Endpoint} & \textbf{Description} \\ \hline
    \multicolumn{3}{|l|}{\textit{Student}} \\ \hline
    GET  & \texttt{/api/students/attendance/summary}  & Attendance overview \\
    GET  & \texttt{/api/students/attendance/sessions} & Session-wise attendance \\
    GET  & \texttt{/api/students/attendance/presence} & Live presence status \\
    POST & \texttt{/api/students/attendance/ping}     & Send attendance ping \\ \hline
    \multicolumn{3}{|l|}{\textit{Admin}} \\ \hline
    GET  & \texttt{/api/admin/attendance/weekly}            & Weekly attendance report \\
    GET  & \texttt{/api/admin/attendance/sessions-by-date}  & Sessions for a date \\
    GET  & \texttt{/api/admin/attendance/session-students}  & Detailed student attendance \\
    GET  & \texttt{/api/admin/attendance/snapshot}          & Snapshot-based attendance \\
    POST & \texttt{/api/admin/attendance/override}          & Manual attendance override \\ \hline
    \multicolumn{3}{|l|}{\textit{Monitoring}} \\ \hline
    GET & \texttt{/api/admin/live-students} & Real-time connected devices \\
    GET & \texttt{/api/admin/wifi-logs}     & Wi-Fi logs for audit \\ \hline
    \multicolumn{3}{|l|}{\textit{Settings}} \\ \hline
    GET    & \texttt{/api/admin/settings/config}         & Get system config \\
    PUT    & \texttt{/api/admin/settings/config}         & Update config \\
    GET    & \texttt{/api/admin/settings/mac-approvals}  & Pending MAC approvals \\
    PATCH  & \texttt{/api/admin/settings/mac-approvals}  & Approve/reject MAC \\
    GET    & \texttt{/api/admin/settings/bssid}          & Get router BSSID \\
    POST   & \texttt{/api/admin/settings/bssid}          & Add BSSID \\
    PATCH  & \texttt{/api/admin/settings/bssid}          & Update BSSID \\
    DELETE & \texttt{/api/admin/settings/bssid}          & Delete BSSID \\ \hline
    \end{tabular}
    \caption{Attendance API endpoints (Student, Admin, Monitoring, and Settings)}
    \label{tab:api_attendance}
\end{table}

\section{Feedback APIs}

\begin{table}[!ht]
    \centering
    \begin{tabular}{|l|l|p{5.5cm}|}
    \hline
    \textbf{Method} & \textbf{Endpoint} & \textbf{Description} \\ \hline
    \multicolumn{3}{|l|}{\textit{Student}} \\ \hline
    GET  & \texttt{/api/feedback/pending}      & Get pending feedback \\
    POST & \texttt{/api/feedback/submit}       & Submit feedback \\
    GET  & \texttt{/api/feedback/my-response}  & View submitted response \\
    GET  & \texttt{/api/feedback/leaderboard}  & Engagement leaderboard \\ \hline
    \multicolumn{3}{|l|}{\textit{Admin}} \\ \hline
    GET    & \texttt{/api/admin/feedback/analytics}       & Feedback analytics \\
    GET    & \texttt{/api/admin/feedback/forms}           & List feedback forms \\
    PATCH  & \texttt{/api/admin/feedback/forms}           & Update deadlines \\
    GET    & \texttt{/api/admin/feedback/questions}       & Get questions \\
    POST   & \texttt{/api/admin/feedback/questions}       & Add question \\
    PATCH  & \texttt{/api/admin/feedback/questions}       & Update question \\
    DELETE & \texttt{/api/admin/feedback/questions}       & Delete question \\
    GET    & \texttt{/api/admin/feedback/student-response}& View individual response \\
    GET    & \texttt{/api/admin/feedback/status}          & Pending feedback status \\ \hline
    \end{tabular}
    \caption{Feedback API endpoints (Student and Admin)}
    \label{tab:api_feedback}
\end{table}

\section{Notification APIs}

\begin{table}[!ht]
    \centering
    \begin{tabular}{|l|l|p{6.5cm}|}
    \hline
    \textbf{Method} & \textbf{Endpoint} & \textbf{Description} \\ \hline
    POST  & \texttt{/api/admin/notifications}   & Send notifications/reminders \\
    GET   & \texttt{/api/student/notifications} & Fetch student notifications \\
    PATCH & \texttt{/api/student/notifications} & Update notification status \\ \hline
    \end{tabular}
    \caption{Notification API endpoints}
    \label{tab:api_notif}
\end{table}

\section{Cron / Automation APIs}

\begin{table}[!ht]
    \centering
    \begin{tabular}{|l|l|p{6.5cm}|}
    \hline
    \textbf{Method} & \textbf{Endpoint} & \textbf{Description} \\ \hline
    GET & \texttt{/api/cron/process-attendance} & Auto process attendance \\
    GET & \texttt{/api/cron/feedback-reminder}  & Send feedback reminders \\ \hline
    \end{tabular}
    \caption{Cron and automation API endpoints}
    \label{tab:api_cron}
\end{table}
